Table~\ref{tab:repair-summary} summarizes the two controlled repairs.
On CIFAR-10, changing the initial covariance at a fixed denoising budget
improves agreement with real images, while a large share of generated
images still falls in a single teacher region. On MNIST, the comparison
starts from the same targeted parent and holds the update budget fixed.
Returning to ordinary training addresses the diagnosed excess and improves
mass balance across digits, with all-class equivalence unresolved.

\begin{table*}[t]
\centering
\caption{Matched-budget outcomes for the two repairs. CIFAR-10 pairs five
fresh generated panels from one checkpoint. The KID interval has nominal
95\% simultaneous coverage across five contrasts under the panel-effect
approximation. MNIST pairs three training
seeds and uses 97.5\% intervals for each of two prespecified outcomes.
Intervals condition on the fixed evaluation references. The coverage row
reports an ordinary paired 95\% interval. The largest teacher region has
reference mass 0.3535. Its row reports pooled masses.}
\label{tab:repair-summary}
\begin{tabular*}{\textwidth}{@{\extracolsep{\fill}}llrrr@{}}
\toprule
Study and matched budget & Outcome & Control & Repair & Paired difference [interval] \\
\midrule
CIFAR-10, two denoiser calls & Kernel Inception Distance $\downarrow$ &
0.1858 & 0.1264 & $-0.05938\;[-0.06138,-0.05739]$ \\
& Real-data coverage $\uparrow$ & 0.206 & 0.288 & $+0.0821\;[0.0672,0.0970]$ \\
& Largest teacher-region mass & 0.9730 & 0.9129 & \textemdash \\
\addlinespace
MNIST, 8,000 updates & Mean digit-mass error $\downarrow$ &
0.02177 & 0.00900 & $-0.01277\;[-0.01878,-0.00676]$ \\
& Digit-3 mass error $\downarrow$ & 0.06837 & 0.00747 &
$-0.06090\;[-0.11814,-0.00366]$ \\
\bottomrule
\end{tabular*}
\end{table*}
